% ============================================================================
% ARBEITSBLATT DS 2 - Gruppe A: Freizeitorte (S. 31-32)
% ============================================================================

\documentclass[11pt,a4paper]{article}

\usepackage[utf8]{inputenc}
\usepackage[T1]{fontenc}
\usepackage[ngerman]{babel}
\usepackage[left=2cm,right=2cm,top=1.8cm,bottom=1.5cm]{geometry}
\usepackage{helvet}
\usepackage[table]{xcolor}
\usepackage{tabularx}
\usepackage{array}
\usepackage{enumitem}
\usepackage{tcolorbox}
\usepackage{fancyhdr}
\usepackage{lastpage}

\renewcommand{\familydefault}{\sfdefault}

\definecolor{spanisch}{HTML}{c41e3a}
\definecolor{grau}{HTML}{666666}
\definecolor{hellgrau}{HTML}{f5f5f5}

\pagestyle{fancy}
\fancyhf{}
\fancyhead[L]{\small\textcolor{grau}{Spanisch Spätbeginner -- Gruppe A}}
\fancyhead[R]{\small\textcolor{grau}{AB-02-A}}
\fancyfoot[C]{\small\thepage/\pageref{LastPage}}
\renewcommand{\headrulewidth}{0.5pt}

\newcommand{\luecke}[1]{\underline{\hspace{#1}}}
\newcommand{\punktebox}[1]{\hfill\fbox{\small \_\_/#1}}
\newcommand{\es}[1]{\textit{#1}}

\newtcolorbox{merkkasten}[1][]{
    colback=white,
    colframe=spanisch,
    fonttitle=\bfseries\color{spanisch},
    title=#1,
    boxrule=1.5pt,
    arc=0pt,
    left=5pt, right=5pt, top=5pt, bottom=5pt
}

\newtcolorbox{buchverweis}{
    colback=hellgrau,
    colframe=grau,
    boxrule=0.5pt,
    arc=0pt,
    left=5pt, right=5pt, top=3pt, bottom=3pt
}

\newcounter{aufgabe}
\newcommand{\aufgabe}[2]{%
    \stepcounter{aufgabe}%
    \vspace{0.8em}%
    \noindent\textbf{\theaufgabe.} #1 \punktebox{#2}%
    \vspace{0.3em}%
}

\begin{document}

% ============================================================================
% KOPFBEREICH
% ============================================================================
\noindent
\begin{tabularx}{\textwidth}{@{}Xr@{}}
    {\Large\textbf{\textcolor{spanisch}{¿Adónde vamos?}}} & Name: \luecke{5cm} \\[0.3em]
    {\small DS 2 -- Freizeitorte (S. 31--32)} & Datum: \luecke{3cm} \\
\end{tabularx}

\vspace{0.3em}
\hrule
\vspace{0.8em}

% ============================================================================
% MERKKASTEN 1: ir a + Ort
% ============================================================================
\begin{merkkasten}[Kasten 1: Das Verb \es{ir} (gehen) + Ort]
\vspace{0.2em}

\begin{tabularx}{\textwidth}{@{}|c|l||c|l|@{}}
\hline
\rowcolor{hellgrau}
\textbf{Person} & \textbf{ir} & \textbf{Person} & \textbf{ir} \\
\hline
yo & \luecke{2cm} & nosotros/-as & \luecke{2cm} \\
\hline
tú & \luecke{2cm} & vosotros/-as & \luecke{2cm} \\
\hline
él/ella/usted & \luecke{2cm} & ellos/-as/ustedes & \luecke{2cm} \\
\hline
\end{tabularx}

\vspace{0.4em}
\textbf{Wichtig:} a + el = \textbf{al} \quad (Beispiel: Voy \textbf{al} cine.)

Aber: a la, a los, a las bleiben getrennt!
\vspace{0.2em}
\end{merkkasten}

\vspace{0.8em}

% ============================================================================
% MERKKASTEN 2: Freizeitorte
% ============================================================================
\begin{merkkasten}[Kasten 2: Freizeitorte -- Lugares de ocio]
\vspace{0.2em}

\begin{tabularx}{\textwidth}{@{}X|X|X@{}}
\es{la piscina} = \luecke{2.5cm} & \es{el cine} = \luecke{2.5cm} & \es{el teatro} = \luecke{2.5cm} \\[0.5em]
\es{el parque} = \luecke{2.5cm} & \es{la discoteca} = \luecke{2.5cm} & \es{el polideportivo} = \luecke{2.5cm} \\[0.5em]
\es{el centro comercial} = \luecke{2.5cm} & \es{la biblioteca} = \luecke{2.5cm} & \es{el estadio} = \luecke{2.5cm} \\
\end{tabularx}

\vspace{0.3em}
\end{merkkasten}

\vspace{1em}

% ============================================================================
% AUFGABEN
% ============================================================================

\aufgabe{Konjugiere das Verb \es{ir}. Ergänze den Kasten oben.}{6}

\vspace{0.3em}

\aufgabe{Ergänze \es{al}, \es{a la}, \es{a los} oder \es{a las}.}{6}

\begin{enumerate}[label=\alph*), leftmargin=2em]
    \item Voy \luecke{1.5cm} piscina todos los sábados.
    \item ¿Vamos \luecke{1.5cm} cine esta noche?
    \item Mis amigos van \luecke{1.5cm} parque después de clase.
    \item Mi hermana va \luecke{1.5cm} biblioteca para estudiar.
    \item Vais \luecke{1.5cm} discotecas los viernes.
    \item Los niños van \luecke{1.5cm} polideportivo.
\end{enumerate}

\vspace{0.5em}

\aufgabe{Hörverstehen: Radio-Spots. (Buch S. 31, Aufg. 2)}{6}

\begin{buchverweis}
\textbf{Hört die Radio-Spots.} Notiert: Welcher Ort? Wann? Was kann man dort machen?
\end{buchverweis}

\vspace{0.3em}
\begin{tabularx}{\textwidth}{@{}|c|X|X|X|@{}}
\hline
\rowcolor{hellgrau}
\textbf{Spot} & \textbf{Ort} & \textbf{Wann?} & \textbf{Aktivität} \\
\hline
1 & \luecke{3cm} & \luecke{3cm} & \luecke{4cm} \\[0.6em]
\hline
2 & \luecke{3cm} & \luecke{3cm} & \luecke{4cm} \\[0.6em]
\hline
3 & \luecke{3cm} & \luecke{3cm} & \luecke{4cm} \\[0.6em]
\hline
\end{tabularx}

\vspace{0.8em}

\aufgabe{Leseverstehen: Martas E-Mail. (Buch S. 32, Aufg. 3)}{8}

\begin{buchverweis}
\textbf{Lies Martas E-Mail.} Beantworte die Fragen auf Spanisch.
\end{buchverweis}

\vspace{0.3em}
\begin{enumerate}[label=\alph*), leftmargin=2em]
    \item ¿Adónde va Marta los sábados?

    \luecke{12cm}

    \item ¿Con quién va al cine?

    \luecke{12cm}

    \item ¿Qué hace en el parque?

    \luecke{12cm}

    \item ¿Por qué le gusta su barrio?

    \luecke{12cm}
\end{enumerate}

\vspace{0.8em}

\aufgabe{Schreibe einen Dialog: Wochenendpläne. (6 Sätze)}{8}

\begin{buchverweis}
\textbf{Situation:} Du und dein Freund plant das Wochenende. Verwendet: \es{¿Adónde vamos?}, \es{Vamos a...}, \es{¿Quieres ir a...?}
\end{buchverweis}

\vspace{0.3em}
\begin{tabularx}{\textwidth}{@{}|X|@{}}
\hline
\\[0.4em]
\textbf{A:} \luecke{13cm} \\[0.7em]
\textbf{B:} \luecke{13cm} \\[0.7em]
\textbf{A:} \luecke{13cm} \\[0.7em]
\textbf{B:} \luecke{13cm} \\[0.7em]
\textbf{A:} \luecke{13cm} \\[0.7em]
\textbf{B:} \luecke{13cm} \\[0.4em]
\hline
\end{tabularx}

\vspace{1em}
\hrule
\vspace{0.3em}
\begin{flushright}
\textbf{Gesamt: \luecke{1cm} / 34 Punkte}
\end{flushright}

\end{document}
